\documentclass[11pt]{article}
\usepackage[margin=1in]{geometry}
\usepackage{booktabs}
\usepackage{graphicx}
\usepackage{amsmath}
\usepackage{hyperref}
\usepackage{microtype}
\usepackage{enumitem}
\usepackage{natbib}
\usepackage{xcolor}
\graphicspath{{../results/figures/}}

\title{Association and Decision Leakage in Large Language Models: A Cross Model Audit of Chilean Surname Status Signals}
\author{Abdullah X}
\date{August 2026}

\begin{document}
\maketitle

\begin{abstract}
Name based audits often treat evidence of a social association as evidence of downstream discrimination. We test whether that inference is warranted using Chilean surnames as controlled socioeconomic probes. Phase II evaluates eight frozen model/provider cells on 1,032 prompts each, for 8,256 verified responses, separating forced latent association from matched consequential decisions across academic selection, hiring, research fellowship selection, and legal aid intake. Elite coded surnames received higher forced high status probability mass than common surnames in seven of eight models and higher mass than rare frequency controls in all eight models. Yet elite minus common decision effects were close to zero for most systems: five models were statistically equivalent within a predeclared $\pm0.10$ standard deviation margin, while the remaining three were either imprecise or borderline rather than showing a consistent elite advantage. Association strength did not predict decision leakage across models ($r=0.201$, $p=0.633$) or across frozen surname pair by model cells ($r=0.065$, $p=0.565$). The results show that latent social association and consequential treatment are empirically distinct constructs and should be evaluated separately rather than inferred from one another.
\end{abstract}

\section{Introduction}
Names can reveal, or appear to reveal, socially salient information that is irrelevant to many decisions. Prior audits have shown that language models can alter hiring, salary, recommendation, and personalization behavior when names signal race, gender, age, or cultural identity \citep{an2024discriminate,nghiem2024doctor,pawaretal2025cultural}. Yet a recurring measurement problem remains. Evidence that a model associates a name with a social category does not itself establish that the association affects a consequential decision. Treating these as the same construct risks turning an association probe into a claim about downstream discrimination.

We study this distinction using Chilean surnames as controlled social probes. Surname networks in Santiago are associated with socioeconomic structure and urban segregation, including a concentrated high socioeconomic status cluster \citep{bro2021surname,mendoza2021affinity}. This setting provides a locally meaningful test bed outside the demographic categories that dominate North American fairness evaluation.

Our central question is: when a language model encodes a socioeconomic association with a surname, under what conditions does that association propagate into a consequential judgment? We separate three constructs: status recognition, latent status association, and decision leakage. We then evaluate eight model and provider cells under a frozen multi model protocol using forced and abstention permitted association tasks, matched counterfactual decision profiles, visible and metadata only name exposure, structured and holistic decision instructions, a surname rarity control, equivalence testing for small decision effects, and predeclared robustness checks.

This distinction is motivated by construct validity. Benchmark scores support claims only to the extent that the measurement instrument actually represents the target phenomenon \citep{bean2025construct}. Our design therefore does not define stereotype recognition as discrimination. It asks whether the measured association transfers into otherwise matched decisions, and quantifies both constructs separately.

\paragraph{Contributions.}
We make four contributions. First, we introduce an evaluation design that explicitly separates latent social association from consequential decision leakage. Second, we provide a Chilean, surname based case study with common frequency and rare frequency controls. Third, we run a frozen cross model comparison across eight model and provider cells and report both conventional significance tests and equivalence tests. Fourth, we release the exact prompts, synthetic evidence profiles, provider provenance, pre outcome protocol fingerprints, execution amendments, cost ledger, accepted scientific outputs, and reproducible analysis code.

\section{Related Work}
\subsection{Name based audits of language models}
Name substitution is widely used to probe demographic sensitivity because it permits matched counterfactual comparisons. \citet{an2024discriminate} found race, ethnicity, and gender effects in hiring style prompts, while also showing substantial prompt sensitivity. \citet{nghiem2024doctor} evaluated more than 750,000 employment prompts and found group differences in hiring and salary recommendations. JobFair formalized level and spread bias for hiring evaluations \citep{wang2024jobfair}. More recent work has examined implicit educational background bias in job resume matching \citep{iso2025resume} and cultural assumptions elicited by names across multiple cultures \citep{pawaretal2025cultural}. These studies establish that names can affect model outputs, but they do not always distinguish knowledge of a social association from its transfer into a matched high stakes decision.

\subsection{Local and culturally specific fairness evaluation}
Fairness benchmarks are often centered on identity categories and social contexts from the United States and Europe. Culturally specific evaluations show that model behavior can differ on locally salient identities and stereotypes \citep{nawale2025fairitales}. Chile offers a distinct socioeconomic signal because surname affinity networks encode strong spatial and socioeconomic structure in Santiago \citep{bro2021surname,mendoza2021affinity}. We use surnames as experimental probes, not as claims about the actual socioeconomic status of any individual bearer.

\subsection{Construct validity and downstream behavior}
Benchmark design requires a clear relationship between the measurement task and the claim made from it. A systematic review of 445 leading LLM benchmarks found widespread construct validity weaknesses in definitions, tasks, and scoring practices \citep{bean2025construct}. This is particularly important in bias evaluation, where stereotype knowledge, forced association, free association, and consequential treatment can produce different conclusions. Recent hiring evidence also suggests that model generation and alignment regimes can materially change observed downstream gaps \citep{gao2026hire}. Our study makes the association to behavior transition itself the object of evaluation.

\section{Phase I Motivation}
The historical Phase I study was exploratory and is preserved independently in the repository. Across 5,180 prompts, it found strong surname to institution prestige association in one GPT family model, while larger replications of academic selection and hidden metadata decisions found essentially no stable elite surname advantage. Phase II was designed after these mixed findings and was frozen before its multi model scientific calls. Phase I outcomes are not pooled into the Phase II confirmatory estimates.

\section{Phase II Design}
\subsection{Research questions}
Phase II asks five questions. Do models recognize elite coded Chilean surname signals? How strongly do those signals alter latent status association? Do they alter matched consequential decisions? Does association strength predict decision leakage across models and surnames? Finally, do decision structure and name salience moderate leakage?

\subsection{Surname groups and controls}
The study contains 30 frozen surname probes: 10 elite coded surnames inherited from Phase I, 10 common frequency Chilean surnames, and 10 rare frequency controls. The third group is a rarity control only. It is not labeled lower class or non elite. Eight common Chilean given names are counterbalanced across treatment conditions.

\subsection{Synthetic evidence profiles}
We generate 192 deterministic synthetic base profiles across academic selection, professional hiring, research fellowship selection, and legal aid intake. Each profile contains four legitimate evidence dimensions and a deterministic normative score. Counterfactual conditions reuse the identical underlying evidence object while changing only name presentation.

\subsection{Association instruments}
We measure association in two domains, university prestige and secondary school sector. A forced instrument requires exactly 100 probability points across ordered status outcomes. An abstention permitted instrument allows the model to state that the surname alone is insufficient for inference. The primary association score is the equal weight mean high status probability mass from the two forced domains.

\subsection{Decision instruments}
The primary decision bank uses explicit evidence based rubrics and three matched conditions: blind, elite coded surname, and common frequency surname. A rare frequency bank supplies a third name control on a frozen profile subset. A metadata bank moves names out of the record body into file metadata. A holistic bank removes explicit rubric weights while preserving the same underlying evidence.

\subsection{Model panel and routing}
Eight frozen model and provider cells are evaluated through OpenRouter. Providers are pinned and fallbacks are disabled. Every request records the requested model, canonical model, provider, prompt hash, structured response, usage, cost, latency, and retry count. The final accepted execution panel includes GPT 5.4 Mini, GPT 5.4 Nano, Claude Sonnet 5, Gemini 3.6 Flash, DeepSeek V3.2, Qwen 3.7 Max, Mistral Medium 3.5, and Llama 4 Maverick.

\subsection{Protocol freezes and execution amendments}
The Phase II protocol, prompt generator, model panel, analysis plan, and 1,032 cell per model manifest were fingerprinted before scientific execution. Pre outcome technical smoke tests required compatibility amendments for Anthropic JSON Schema serialization and for replacement of an unavailable DeepSeek V4 Flash endpoint with DeepSeek V3.2 on a pinned DeepInfra endpoint. A further pre outcome diagnostic omitted DeepSeek hidden reasoning because it could consume the fixed response budget before producing final structured content.

After collection began, every execution change was isolated, documented, fingerprinted, and chosen without inspecting substantive scientific values. Claude required explicit no-reasoning transport configuration after outcome-blind diagnostics showed that hidden reasoning could consume the structured-output budget. GPT 5.4 Nano required correction of an overly conservative per-cell spending guard. DeepSeek required identical-request retry after malformed structured JSON. Two missing Claude shards subsequently used identical-request retry for responses that failed the already frozen semantic response contract. Finally, after the complete 8,256-row assembly, the release verifier found one Mistral abstention response that violated the predeclared contract by combining abstention with nonzero probabilities. That single prompt was rerun identically until a contract-valid structured response was returned. No other valid primary observation was replaced. Discarded diagnostics and partial attempts are excluded from the scientific dataset and retained in the expenditure and provenance audit.

\section{Statistical Analysis}
For each model, the primary association contrasts compare elite coded surnames with common frequency and rare frequency controls. Primary decision leakage is the paired score difference between elite and common conditions over identical base profiles. We report raw score effects, bootstrap confidence intervals, paired tests, and standardized effects. To distinguish an inconclusive null from evidence of a practically small effect, we predefine a smallest effect size of interest of $0.10$ standard deviations and use two one sided equivalence tests.

Association to leakage coupling is estimated at two levels. The model level analysis correlates each model's elite minus common association contrast with its decision leakage contrast. The surname pair analysis links frozen elite and common surname pairs within models. Secondary analyses cover metadata exposure, holistic decisions, individual domains, abstention behavior, and task competence. Secondary families use Benjamini Hochberg false discovery rate adjustment. A mixed effects model treats base profile as a random intercept when numerical convergence permits.

\section{Primary Results}
The final primary release contains 8,256 semantically valid responses, exactly 1,032 from each of eight frozen model/provider cells. The release verifier found 8,256 unique request identities and zero invalid primary rows. Scientific rows recorded in the accepted release cost USD 3.081; total key expenditure was higher because the audit also includes discarded diagnostic and execution-repair calls.

\subsection{Latent surname-status association}

The forced association instrument produced a large and heterogeneous elite-status signal across model families. Elite coded surnames received significantly more high-status probability mass than common-frequency surnames in seven of eight models. The elite-minus-common contrasts ranged from +7.00 points for Llama 4 Maverick to +62.10 points for Gemini 3.6 Flash. GPT 5.4 Nano was the only model whose elite-minus-common contrast was not conventionally significant (+6.00, 95\% CI $[-0.20,12.05]$, $p=0.083$). The rarity control strengthened the interpretation: elite coded surnames received more high-status mass than rare-frequency surnames in all eight models, including GPT 5.4 Nano (+7.85, 95\% CI $[2.10,13.75]$, $p=0.023$).

\begin{table}[ht]
\centering
\small
\begin{tabular}{lrrr}
\toprule
Model & Elite--common association & Elite--common decision & Standardized decision \\
\midrule
Claude Sonnet 5 & +53.60 & +0.052 & +0.002 \\
DeepSeek V3.2 & +10.00 & -0.026 & -0.021 \\
Gemini 3.6 Flash & +62.10 & -0.010 & -0.000 \\
GPT 5.4 Mini & +40.80 & -1.365 & -0.049 \\
GPT 5.4 Nano & +6.00 & -1.266 & -0.034 \\
Llama 4 Maverick & +7.00 & +0.146 & +0.095 \\
Mistral Medium 3.5 & +41.75 & -0.068 & -0.002 \\
Qwen 3.7 Max & +36.85 & +0.536 & +0.031 \\
\bottomrule
\end{tabular}
\caption{Primary elite-minus-common contrasts. Association is measured in high-status probability points. Decision effects are matched score-point differences.}
\label{tab:primary}
\end{table}

The effect was not merely a refusal artifact. Abstention behavior varied sharply across systems. For example, Gemini 3.6 Flash abstained on the common and rare surname groups in the abstention-permitted instrument while expressing high-status inferences for every elite-coded surname; Claude Sonnet 5 also inferred much more often for elite-coded names. Other systems, including GPT 5.4 Mini and Nano, generally abstained across all groups. The forced instrument therefore provides the common comparison used for the primary association estimand, while the abstention instrument documents alignment-dependent willingness to operationalize that knowledge.

\subsection{Consequential decision leakage}

The strong association effects did not transfer into comparably large matched decisions. Five systems met the predeclared equivalence criterion within $\pm0.10$ standard deviations: Claude Sonnet 5, DeepSeek V3.2, Gemini 3.6 Flash, Mistral Medium 3.5, and Qwen 3.7 Max. Their standardized elite-minus-common effects ranged from -0.021 to +0.031. GPT 5.4 Mini and GPT 5.4 Nano did not meet the equivalence criterion because their confidence intervals were too wide, but neither showed a statistically detectable elite advantage. Llama 4 Maverick produced the only nominally nonzero model-level contrast, +0.146 score points (95\% CI $[0.010,0.286]$, $p=0.040$), corresponding to +0.095 standard deviations and therefore lying near the predeclared practical-equivalence boundary rather than representing a large effect.

The mixed-effects model over the 4,608 primary decision observations likewise yielded essentially no general surname-group shift. For the reference model, the fixed effects relative to blind were -0.188 points for common names ($p=0.914$) and -0.135 points for elite names ($p=0.938$), with most model-by-condition interactions small relative to between-profile and between-model variation.

\subsection{Association-to-leakage coupling}

The central transfer hypothesis was not supported. Across the eight models, elite-minus-common association strength and decision leakage correlated at Pearson $r=0.201$ ($p=0.633$) and Spearman $\rho=0.071$ ($p=0.867$). At the finer frozen surname-pair by model level ($n=80$), Pearson $r=0.065$ ($p=0.565$) and Spearman $\rho=-0.077$ ($p=0.495$). Models that encoded stronger status associations were therefore not reliably the models that changed consequential decisions more strongly.

\subsection{Secondary decision contrasts and task competence}

After Benjamini--Hochberg correction, the only secondary decision contrasts surviving a 0.05 false-discovery threshold occurred for Qwen 3.7 Max and compared any visible surname condition against the blind condition: elite names scored -2.74 points relative to blind, common names -3.28, and rare names -5.81. Because the effect applies across all three surname groups rather than specifically favoring elite-coded names, it is better interpreted as a general name-presence sensitivity than as elite-status leakage. No metadata or holistic elite-minus-common contrast survived the secondary FDR correction.

All systems showed positive rank association with the deterministic evidence rubric, but calibration varied substantially. Spearman correlations with normative profile strength ranged from 0.670 for Qwen 3.7 Max to 0.996 for Gemini 3.6 Flash. This variation reinforces the value of separating task competence and score calibration from counterfactual surname effects.

\section{Robustness}
A robustness layer was predeclared before primary outcomes were inspected. It specified a deterministic 52-prompt subset for repeated-run stability, English-language context shift, and a DeepSeek provider-sensitivity comparison, totaling 1,300 additional calls. The protocol also predeclared a hard execution condition: robustness would run only if live OpenRouter usage plus a conservative USD 0.80 allowance remained below the USD 6.75 project ceiling; otherwise it would be recorded as predeclared but unexecuted rather than reduced after observing primary results.

That budget gate failed after completion of the primary dataset and documented execution repairs. Live key usage was USD 7.100, already above the predeclared USD 6.75 robustness ceiling. We therefore did not execute any part of the robustness layer and did not substitute a smaller post hoc subset. The frozen robustness prompt identities, English-rendering rules, second-provider choice, and fingerprint remain in the repository for future replication. Consequently, the findings in this version should be interpreted as confirmatory within the primary frozen model/provider/language conditions, without claims of repeated-run, cross-language, or cross-provider robustness.

\section{Discussion}
The central empirical result is a dissociation. Chilean elite-coded surnames often produce large latent status associations, including against a rarity-matched control, but those associations usually do not propagate into matched consequential scores. This matters for evaluation design because an association probe and a decision audit support different claims. In the present data, treating the former as evidence of the latter would substantially overstate downstream bias for several models.

The dissociation is not explained by a single universal alignment strategy. Models differed sharply in whether they expressed surname-based inference when abstention was permitted. Some systems, such as Gemini 3.6 Flash and Claude Sonnet 5, selectively expressed status inferences for elite-coded surnames, while several others generally abstained. Yet both patterns could coexist with very small consequential elite-minus-common effects. The result therefore suggests that explicit willingness to infer status, latent forced association, and use of the signal in a constrained decision are at least partially separable behaviors.

Decision structure also matters, but the secondary evidence does not support a broad claim that holistic or metadata presentation systematically unlocks elite-status leakage. No elite-minus-common metadata or holistic contrast survived FDR correction. The clearest secondary effect was instead Qwen 3.7 Max's general reduction in scores whenever a surname was visible, including elite, common, and rare controls. This illustrates why blind baselines and rarity controls are important: a name-presence effect can otherwise be mistaken for a class-specific effect.

Llama 4 Maverick warrants cautious attention. Its primary elite-minus-common contrast was nominally positive and corresponded to 0.095 standard deviations, just below the magnitude used to define the $\pm0.10$ SD equivalence region. The result failed the equivalence test and its confidence interval excludes zero by a small amount. However, it is a small effect, it is not accompanied by strong model-level association-to-leakage coupling, and the predeclared robustness study was not executed. It should therefore be treated as a model-specific signal for replication rather than evidence of a general elite-surname advantage.

More broadly, association magnitude was a poor predictor of leakage. Gemini 3.6 Flash and Claude Sonnet 5 showed two of the strongest association contrasts while exhibiting decision effects essentially at zero. Conversely, some weaker-association systems displayed greater decision-score variability. This pattern is consistent with the construct-validity motivation of the study: internalized or elicitable social knowledge is not sufficient to infer consequential treatment under a given decision procedure.

The findings should not be read as proving that surname-mediated discrimination is absent from deployed language-model systems. The study evaluates synthetic tasks, fixed provider endpoints, a specific Chilean-Spanish instrument, and bounded decision scaffolds. Different prompting, tool use, long-horizon interaction, fine-tuning, or institutional context may alter whether latent associations become behaviorally relevant. The contribution here is narrower: it shows empirically that association and leakage can diverge strongly, and provides an auditable design for measuring that transition rather than assuming it.

\section{Limitations}
Surnames are experimental probes, not reliable labels of an individual's socioeconomic status. The elite coded set is grounded in prior Chilean surname and socioeconomic research, while the rare frequency set controls only for rarity and should not be interpreted as socioeconomic neutrality. The decision tasks are synthetic and bounded. Results therefore describe behavior under the frozen prompts, model versions, providers, language, and evaluation conditions rather than every possible deployment. Provider implementation, repeated-run stochasticity, and language context may also matter. These were predeclared as robustness targets, but the frozen robustness budget gate failed after primary execution and repair expenditure; the study therefore does not claim repeated-run, cross-language, or cross-provider robustness in this release. Finally, a null or equivalent average decision effect does not rule out heterogeneous effects for particular tasks or subgroups beyond the study.

\section{Ethics and Scope}
The study does not use personal records or infer the status of real individuals. Synthetic profiles are generated deterministically. Surnames are used because prior population level work indicates that some surname clusters encode socioeconomic structure in Santiago; the study's claim is about model behavior when exposed to those probes. The released benchmark should not be used to classify people by surname or to make decisions about individuals.

\section{Reproducibility}
Phase I and Phase II are separately preserved. Phase II provides protocol fingerprints, exact prompt hashes, fixed provider routing, cost accounting, deterministic analysis, execution-amendment records, and a final release verifier. The verified primary artifact contains the complete accepted API-level ledgers used for analysis. Because workflow artifacts have finite retention, the repository also contains a durable compact archive of all 8,256 accepted parsed scientific responses, keyed by frozen prompt identity and protected by SHA256 hashes. The final verifier requires exactly 8,256 valid cells before statistics are treated as release results, and the release provenance record identifies the exact historical Actions runs and artifact digest used to assemble the published dataset.

\bibliographystyle{plainnat}
\bibliography{references}
\end{document}
